

Explicación #2: Señal territorial (interpretación exploratoria)

Estas reglas capturan **señales territoriales** moderadas. Son útiles para describir diferencias en el registro, pero no deben interpretarse causalmente (pueden mezclar incidencia real + prácticas administrativas + calidad de captura).

\begin{table}[htbp]
\centering
\small
\begin{tabularx}{\linewidth}{>{\raggedright\arraybackslash}p{1.2cm} X X >{\raggedleft\arraybackslash}p{1.6cm} >{\raggedleft\arraybackslash}p{1.8cm} >{\raggedleft\arraybackslash}p{1.4cm}}
\toprule
ID & Antecedente & Consecuente & Support & Conf. & Lift \\
\midrule
22 & \texttt{ENTIDAD=ESTADO DE MÉXICO} & \texttt{ESTATUS\_VICTIMA=DESAPARECIDA} & 0.0775 & 0.7235 & 1.213 \\
\bottomrule
\end{tabularx}
\caption{Reglas territoriales (lectura exploratoria). Valores redondeados.}
\label{tab:territorial}
\end{table}

Esto se debe a que…** un `lift` alrededor de `1.21` sugiere una asociación por encima del promedio del dataset. La utilidad aquí es **comparativa/exploratoria**: indica que en esa entidad el estatus “DESAPARECIDA” aparece relativamente más que lo esperado, pero puede estar influido por **criterios de clasificación** y **rezagos de actualización**.

Explicación #3: Sesgo del registro — “confidencialidad en bloque” (censura/datos faltantes no aleatorios)

Estas reglas muestran un patrón sistémico: cuando un campo clave es `CONFIDENCIAL`, otros campos también lo son con **confianza 1.0**, en una porción muy grande del dataset.

\begin{table}[htbp]
\centering
\small
\begin{tabularx}{\linewidth}{>{\raggedright\arraybackslash}p{1.2cm} X X >{\raggedleft\arraybackslash}p{1.6cm} >{\raggedleft\arraybackslash}p{1.8cm} >{\raggedleft\arraybackslash}p{1.4cm}}
\toprule
ID & Antecedente & Consecuente & Support & Conf. & Lift \\
\midrule
7  & \texttt{ESTATUS\_VICTIMA=CONFIDENCIAL} & \texttt{SEXO=CONFIDENCIAL} & 0.3671 & 1.0000 & 2.724 \\
10 & \texttt{SEXO=CONFIDENCIAL} & \texttt{MUNICIPIO=CONFIDENCIAL} & 0.3671 & 1.0000 & 2.724 \\
20 & \texttt{ESTATUS\_VICTIMA=CONFIDENCIAL} & \texttt{MUNICIPIO=CONFIDENCIAL} & 0.3671 & 1.0000 & 2.724 \\
\bottomrule
\end{tabularx}
\caption{Sesgo de registro: confidencialidad en bloque (censura/datos faltantes estructurados). Valores redondeados.}
\label{tab:confidencialidad_bloque}
\end{table}

Esto se debe a que…** con `confidence=1.0` y `support≈0.367`, no estamos ante “una correlación social”, sino ante un **mecanismo de captura/anonimización**: el sistema (o la fuente) tiende a marcar **múltiples variables como confidenciales al mismo tiempo**. Esto introduce un sesgo fuerte: los registros con información completa no son una muestra aleatoria del total.

Explicación #4: Sesgo por entidad — caso Jalisco (diferencias locales en resguardo/captura)

Estas reglas muestran que en `ENTIDAD=JALISCO` hay una asociación marcada con campos `CONFIDENCIAL` y, además, que dentro de Jalisco la confidencialidad aparece “pegada” entre variables.

\begin{table}[htbp]
\centering
\small
\begin{tabularx}{\linewidth}{>{\raggedright\arraybackslash}p{1.2cm} X X >{\raggedleft\arraybackslash}p{1.6cm} >{\raggedleft\arraybackslash}p{1.8cm} >{\raggedleft\arraybackslash}p{1.4cm}}
\toprule
ID & Antecedente & Consecuente & Support & Conf. & Lift \\
\midrule
9  & \texttt{ENTIDAD=JALISCO} & \texttt{SEXO=CONFIDENCIAL} & 0.0747 & 0.6752 & 1.839 \\
19 & \texttt{ENTIDAD=JALISCO} & \texttt{ESTATUS\_VICTIMA=CONFIDENCIAL} & 0.0747 & 0.6752 & 1.839 \\
30 & \texttt{ENTIDAD=JALISCO} & \texttt{MUNICIPIO=CONFIDENCIAL} & 0.0747 & 0.6752 & 1.839 \\
36 & \texttt{ESTATUS\_VICTIMA=CONFIDENCIAL, ENTIDAD=JALISCO} & \texttt{SEXO=CONFIDENCIAL} & 0.0747 & 1.0000 & 2.724 \\
45 & \texttt{ENTIDAD=JALISCO', 'SEXO=CONFIDENCIAL} & \texttt{MUNICIPIO=CONFIDENCIAL} & 0.0747 & 1.0000 & 2.724 \\
46 & \texttt{ENTIDAD=JALISCO', 'MUNICIPIO=CONFIDENCIAL} & \texttt{SEXO=CONFIDENCIAL} & 0.0747 & 1.0000 & 2.724 \\
\bottomrule
\end{tabularx}
\caption{Sesgo por entidad: patrones de confidencialidad asociados a Jalisco. Valores redondeados.}
\label{tab:jalisco_confidencialidad}
\end{table}

Esto se debe a que…** el `confidence≈0.675` desde `ENTIDAD=JALISCO` hacia `*=CONFIDENCIAL` sugiere que una fracción muy grande de registros en esa entidad cae en el “modo confidencial”. El bloque con `confidence=1.0` indica que, **cuando entra confidencialidad**, tiende a hacerlo en paquete. En términos de sesgo, esto dificulta comparar Jalisco con otras entidades en variables como sexo/municipio/estatus sin controlar esta censura.

Explicación #5: Sesgo temporal — rezago de actualización (right-censoring) en años recientes

Estas reglas muestran que para 2024 y 2025 casi todo queda etiquetado como `DESAPARECIDA`, lo cual es consistente con rezagos (casos recientes aún no se actualizan a otros estatus).

\begin{table}[htbp]
\centering
\small
\begin{tabularx}{\linewidth}{>{\raggedright\arraybackslash}p{1.2cm} X X >{\raggedleft\arraybackslash}p{1.6cm} >{\raggedleft\arraybackslash}p{1.8cm} >{\raggedleft\arraybackslash}p{1.4cm}}
\toprule
ID & Antecedente & Consecuente & Support & Conf. & Lift \\
\midrule
28 & \texttt{AÑO\_DESAPARICION=2025.0} & \texttt{ESTATUS\_VICTIMA=DESAPARECIDA} & 0.0512 & 0.9985 & 1.674 \\
29 & \texttt{AÑO\_DESAPARICION=2024.0} & \texttt{ESTATUS\_VICTIMA=DESAPARECIDA} & 0.0598 & 0.9960 & 1.669 \\
\bottomrule
\end{tabularx}
\caption{Sesgo temporal: asociación entre año reciente y estatus ``DESAPARECIDA'' (posible rezago de actualización). Valores redondeados.}
\label{tab:sesgo_temporal}
\end{table}

Esto se debe a que w`confidence≈1.0` en años recientes es una señal típica de **censura por la derecha**: el estatus es “transitorio” y la base no está igualmente actualizada para todos los años. Por eso, cualquier comparación temporal de estatus debe hacerse controlando este efecto (o excluyendo años recientes / usando ventanas de actualización).